%% Quinn Booth: resume
%%
%% Build:    latexmk          (or just save in VS Code with LaTeX Workshop)
%% Preview:  latexmk -pvc     (rebuilds and refreshes on every save)
%% Output:   build/resume.pdf
%%
%% Pushing to main publishes it to
%% https://quinnbooth.github.io/resumes/Quinn_Booth_Resume.pdf
%%
%% Uses only packages every TeX install ships, so it builds the same locally
%% (TinyTeX) and on GitHub (full TeX Live).

\documentclass[11pt, letterpaper]{article}

\usepackage[margin=0.55in, top=0.4in, bottom=0.4in]{geometry}
\usepackage[T1]{fontenc}
\usepackage[utf8]{inputenc}
\renewcommand{\rmdefault}{ptm}  % Times
\usepackage{xcolor}
\usepackage[hidelinks]{hyperref}

% Links worth clicking (papers, press) get a quiet ink blue so they read as
% links without looking like a web page. Header links stay black.
\definecolor{linkink}{HTML}{1F3F73}
\newcommand{\link}[2]{\href{#1}{\textcolor{linkink}{#2}}}

% Real text in the PDF, so copy-paste and applicant-tracking systems read it.
\input{glyphtounicode}
\pdfgentounicode=1

\hypersetup{
  pdftitle={Quinn Booth - Resume},
  pdfauthor={Quinn Booth},
}

\pagestyle{empty}
\setlength{\parindent}{0pt}
\raggedright
\frenchspacing

%% ---- Building blocks ------------------------------------------------------

% A section heading with a thin rule under it.
\newcommand{\sectionhead}[1]{%
  \par\vspace{4pt}%
  {\large\scshape #1}\par
  \vspace{-6pt}\rule{\linewidth}{0.4pt}\par
  \vspace{1pt}}

% \role{Organization}{Place}{Title}{Dates}
\newcommand{\role}[4]{%
  \par\vspace{3pt}%
  \textbf{#1} \hfill #2\par
  \textit{#3} \hfill #4\par}

% \subrole{Title}{Dates}: a second title under the same organization.
\newcommand{\subrole}[2]{%
  \par\vspace{2pt}%
  \textit{#1} \hfill #2\par}

% Tight bullet list.
\newenvironment{bullets}{%
  \begin{list}{\textbullet}{%
    \setlength{\leftmargin}{1.1em}%
    \setlength{\labelsep}{0.45em}%
    \setlength{\labelwidth}{0.65em}%
    \setlength{\itemsep}{0.5pt}%
    \setlength{\parsep}{0pt}%
    \setlength{\topsep}{1.5pt}%
    \setlength{\partopsep}{0pt}}}%
  {\end{list}}

% Smaller text for asides: tech lists, advisors, the YC batch.
\newcommand{\meta}[1]{{\small #1}}

%% ---- Resume ---------------------------------------------------------------

\begin{document}

\begin{center}
  {\Huge Quinn Booth}\par
  \vspace{5pt}
  \small
  San Francisco, CA
  \enspace$\cdot$\enspace \href{mailto:qabooth@gmail.com}{qabooth@gmail.com}
  \enspace$\cdot$\enspace \href{https://quinnbooth.github.io}{quinnbooth.github.io}
  \enspace$\cdot$\enspace \href{https://github.com/quinnbooth}{github.com/quinnbooth}
  \enspace$\cdot$\enspace \href{https://www.linkedin.com/in/quinn-booth}{linkedin.com/in/quinn-booth}
\end{center}

%% ---------------------------------------------------------------------------
\sectionhead{Experience}

\role{Anchor Space}{San Francisco, CA}{Co-founder \meta{(Y Combinator W27)}}{2026 -- Present}
\begin{bullets}
  \item Building low-cost, software-defined laser ground stations that downlink 10x more satellite data than radio.
\end{bullets}

\role{Astranis}{San Francisco, CA}{Flight Software Engineer II}{Dec 2025 -- Nov 2026}
\subrole{Flight Software Engineer}{Sep 2024 -- Nov 2025}
\begin{bullets}
  \item Led and integrated secure boot for our satellites' flight computers, a chain of trust from boot ROM through Arm
        Trusted Firmware (TF-A) and U-Boot verified boot up to flight software. Slated for 5 satellites, dozens more to follow.
  \item Implemented the secure boot PKI, from HSM-backed key generation and RSA-signed FIT images to eFuse provisioning
        that permanently burns the root of trust into each flight computer.
  \item Designed the secure release pipeline, where isolated runners with locked-down ingress and egress verify each release
        is a trusted secure build before signing it in a cloud HSM. Added a dedicated Bazel cache to speed up builds.
  \item Re-architected the hosted payload system, which runs partners' software in QEMU virtual machines on our flight
        computers with a minimal Yocto guest image, CPU pinning, memory caps, and automatic restarts. Wrote the quarantined
        protobuf API that lets guests reach flight services only through allowlists and rate limits.
  \item Launched live telemetry monitoring, with rules that flag off-nominal spacecraft behavior and alert the team.
        Co-developed it with an intern I hired and mentored, who earned a return offer.
  \item Built an automated data review platform that queries telemetry and generates reports on a schedule in GitHub
        Actions. Adopted by 10+ teams, many outside software, for dozens of recurring reports.
  \item Cut on-call pages by up to 50\% by automating triage of radiation-induced upsets (SEUs), with new telemetry,
        Grafana dashboards, reboot procedures, and warm boot on the flight computer. Trained operators to handle them.
  \item Owned a major flight software release end to end, from testbed verification and bug triage to deployment in orbit.
\end{bullets}

%% ---------------------------------------------------------------------------
\sectionhead{Research}

\role{Creative Machines Lab, Columbia University}{New York, NY}{Hardware Lead \& Research Assistant \meta{(Prof.\ Hod Lipson)}}{May 2022 -- May 2024}
\begin{bullets}
  \item Led hardware for Robot Metabolism, truss robots that grow by absorbing parts from other robots, and worked on
        the computer vision and cloud control that ran them. Featured in a \link{https://www.youtube.com/watch?v=sK5_pQV_QEA}{Bloomberg Originals documentary}.
  \item Developed cloud control (C++, REST) for a fleet of 16+ vibrating Particle Robots and designed their PCBs.
\end{bullets}

\role{Internet Real-Time Lab, Columbia University}{New York, NY}{Research Assistant \meta{(Prof.\ Henning Schulzrinne)}}{Sep 2023 -- Aug 2024}
\begin{bullets}
  \item Wrote C++ firmware for low-power LoRa IoT devices, improving energy efficiency more than 10x.
\end{bullets}

%% ---------------------------------------------------------------------------
\sectionhead{Publications}

\begin{bullets}
  \item \link{https://www.science.org/doi/10.1126/sciadv.adu6897}{Robot metabolism: Toward machines that can grow by consuming other machines.}
        \textit{Science Advances}, 2025.
  \item \link{https://dl.acm.org/doi/10.1145/3803784.3816836}{Clay for the Creative Mind: Hands-On Robot Shaping Using Truss Links.}
        \textit{ACM Creativity \& Cognition}, 2026.
  \item \link{https://ieeexplore.ieee.org/document/10619984}{Robot Links: Towards Self-Assembling Truss Robots.}
        \textit{IEEE ReMAR}, 2024.
  \item \link{https://ieeexplore.ieee.org/document/10619983}{Vibrating Particle Robot.}
        \textit{IEEE ReMAR}, 2024.
\end{bullets}

%% ---------------------------------------------------------------------------
\sectionhead{Education}

\par\vspace{4pt}%
\textbf{Columbia University} \enspace B.S. Electrical Engineering, minor in Computer Science
\meta{(GPA 3.94, Dean's List)} \hfill 2020 -- 2024\par

%% ---------------------------------------------------------------------------
\sectionhead{Skills}

\textbf{Languages:} C++, Python, C, SystemVerilog, Verilog, SQL, JavaScript\par
\textbf{Security:} Secure boot, PKI, HSMs, eFuse provisioning, TF-A, U-Boot, FIT image signing, RSA\par
\textbf{Systems:} Flight software, Yocto, QEMU, protobuf, Bazel, pytest, Grafana, GCP, AWS, Git, FPGAs, PCB design

\end{document}
